\documentclass[aspectratio=169]{beamer}
\usepackage[T1]{fontenc}
\usepackage{mathpazo}
\usepackage{tgpagella}
\usepackage{tgheros}
\usepackage{amsmath}
\usefonttheme{professionalfonts}
\providecommand{\DerivationRepositoryRoot}{.}
\providecommand{\DerivationPalettePrefix}{toaster}
\providecommand{\DerivationAssets}{animations/generated}
\input{\DerivationRepositoryRoot/tools/vendor/manim_pdfpc/derivation_toolkit/tex/theme.tex}
\input{\DerivationRepositoryRoot/tools/vendor/manim_pdfpc/derivation_toolkit/tex/derivations.tex}
\setbeamertemplate{footline}[frame number]
\title{Pop-time speedups}
\subtitle{The geometric mean in the fictional toaster report}
\author{Illustrative derivation}
\date{}
\begin{document}
\begin{frame}
  \titlepage
\end{frame}
\begin{frame}{From trial ratios to one speedup}
  \DerivationScene{geometric-mean}
\end{frame}
\begin{frame}{Read the result}
  For the illustrative trial ratios, the equal per-trial factor is
  \[ \DerivationResult{toaster-example} \]
  For positive trial times, averaging the log ratios gives the same result.
\end{frame}
\end{document}
